\section{An arithmetic torus with dense cyclic generators}\label{sec:arithmetic}
The following elementary unit construction fixes the number field before
any genericity condition is imposed.

\begin{lemma}\label{lem:units}
For every $r\geq2$ there are a number field $L\subset\C$ and units
$\mu_1,\ldots,\mu_r\in\OO_L^\times$ with product one such that
\[
 \prod_i\mu_i^{e_i}=1\quad\Longleftrightarrow\quad
 e_1=\cdots=e_r\qquad(e_i\in\Z).
\]
In particular, $C_\mu=\diag(\mu_1,\ldots,\mu_r)$ generates a Zariski-dense
cyclic subgroup of $T$ and has distinct eigenvalues.
\end{lemma}
\begin{proof}
Fix an integer $c_0\geq3$ and set $f(X)=X^r-c_0X^{r-1}-1$.
On $|X|=1$, the middle term has modulus $c_0>2\geq|X^r-1|$.
Rouch\'e's theorem gives exactly $r-1$ roots inside the unit circle.
Since $f(c_0)=-1$ and $f(x)$ tends to infinity on the positive real axis,
there is one real root greater than $c_0$, the unique root outside.
If $f$ were reducible over $\Q$, Gauss's lemma would give a nonconstant
monic integral factor not containing that outside root. Its constant term
would be a nonzero integer of modulus less than one, a contradiction.
Thus $f$ is irreducible; characteristic zero gives distinct roots
$\alpha_1,\ldots,\alpha_r$. They are algebraic units because $f$ is monic
with constant term $-1$.

Suppose $\prod_i\alpha_i^{e_i}$ is a root of unity. Choose an index whose
exponent is $M=\max_i e_i$. Irreducibility allows a Galois automorphism
$\sigma$ of the splitting field that sends this root to the unique outside
root. Since $|\prod_i\sigma\alpha_i|=1$, we have
\[
 0=\sum_i e_i\log|\sigma\alpha_i|
   =\sum_i(e_i-M)\log|\sigma\alpha_i|.
\]
The term for the outside root is zero; all other logarithms are strictly
negative and all their coefficients are nonpositive. Every term is
therefore nonnegative, and equality forces $e_i=M$ for every $i$.

Choose a root of unity $\eta$ with $\eta^r=(-1)^{r-1}$, and put
$\mu_i=\eta\alpha_i$. Since $\prod_i\alpha_i=(-1)^{r+1}$, their product is
one. A relation among the $\mu_i$ gives a root-of-unity relation among the
$\alpha_i$, so its exponents are all equal. Conversely all equal exponents
give a relation. Take $L$ to be the splitting field with $\eta$ adjoined.
The character lattice of $T$ is $\Z^r/\Z(1,\ldots,1)$. A proper algebraic
subgroup of a torus is annihilated by a nontrivial character, so the stated
relation lattice proves cyclic density. Distinctness follows already from
that of the $\alpha_i$.
\end{proof}

For a tuple $m=(m_1,\ldots,m_{r-1})$ of positive integers define
\begin{equation}\label{eq:Bfamily}
 B(m)=\diag\left(\mu_1^{m_1},\ldots,\mu_{r-1}^{m_{r-1}},
                    \prod_{i<r}\mu_i^{-m_i}\right).
\end{equation}
\begin{lemma}\label{lem:arithmetic-dense}
The set of pairs $(B(m),Q)$, with all $m_i\geq1$ and
$Q\in\SL_r(\Z)$, is Zariski dense in $T\times\SL_r(\C)$.
Every $B(m)$ generates $T$ densely. Every
$C=QC_\mu Q^{-1}$ belongs to $\SL_r(\OO_L)$ and generates $\TC$ densely.
\end{lemma}
\begin{proof}
Each $\mu_i$ has infinite order, by Lemma~\ref{lem:units}. The first $r-1$
coordinates in \eqref{eq:Bfamily} range independently over infinite subsets
of $\C^\times$. A Laurent polynomial vanishing on their product is zero,
by induction on the number of variables. This proves density of the
$B(m)$. If a character with exponents $e_i$ is one on $B(m)$, then
\[
 \prod_{i<r}\mu_i^{m_i(e_i-e_r)}=1.
\]
Apply Lemma~\ref{lem:units} with final exponent zero. All exponents are
zero, hence all $e_i=e_r$, proving cyclic density.

The Zariski closure of $\SL_r(\Z)$ contains each elementary root group:
it contains its integral parameters, an infinite subset of an affine line.
Those root groups generate $\SL_r(\C)$, proving density. Products of dense
sets are dense. Finally $Q^{-1}$ has integral entries because $\det Q=1$,
so $QC_\mu Q^{-1}$ is integral; its torus density follows by conjugation.
\end{proof}

All later parameter opens will be taken inside this single fixed-eigenvalue
family. We never assume that a countable set of arithmetic points avoids a
countable union of proper closed sets.
